\section{Related Work}

\subsection{CNN-Based Medical Image Segmentation}

U-Net \cite{ronneberger2015unet} established the encoder-decoder architecture with skip connections as the foundation for medical image segmentation. The symmetric structure captures multi-scale features while skip connections preserve spatial details lost during downsampling. However, U-Net uses max pooling for downsampling, which discards spatial information irreversibly, and batch normalization, which assumes large batch sizes impractical for medical imaging with high-resolution 3D volumes.

Subsequent works addressed these limitations. Attention U-Net \cite{oktay2018attention} introduced attention gates to suppress irrelevant features in skip connections, improving focus on target structures. U-Net++ \cite{zhou2018unet++} employed nested dense skip connections for better gradient flow. nnU-Net \cite{isensee2021nnu} achieved state-of-the-art results through careful hyperparameter tuning and training strategies, including instance normalization and LeakyReLU activation. **In contrast to these pure CNN architectures**, our AMT-UNet incorporates a transformer bottleneck for explicit global context modeling while adopting best practices from nnU-Net (instance normalization, LeakyReLU, residual connections). Additionally, we replace max pooling with learned strided convolutions and enhance skip connections with CBAM attention rather than simple gating.

\subsection{Transformer-Based Medical Image Segmentation}

The success of Vision Transformers (ViT) \cite{dosovitskiy2020image} in natural image recognition inspired their application to medical imaging. TransUNet \cite{chen2021transunet} combined CNN encoders with transformer layers, showing improved performance through global context modeling. UNETR \cite{hatamizadeh2022unetr} used pure transformer encoders with CNN decoders for 3D medical image segmentation. Swin-Unet \cite{cao2022swin} employed hierarchical shifted window attention for multi-scale feature learning with reduced computational cost.

While these methods demonstrate the value of transformers for capturing long-range dependencies, they face limitations in medical imaging: (1) standard self-attention has $O(N^2)$ complexity with respect to sequence length, making it expensive for high-resolution images; (2) transformers require large datasets for effective training, which are scarce in medical domains; (3) pure transformer architectures may lose fine-grained spatial details critical for precise boundary delineation. **AMT-UNet addresses these issues** through a hybrid design that uses CNNs for spatial feature extraction and transformers only at the bottleneck where spatial resolution is lowest. Crucially, our adaptive masking mechanism learns to focus transformer computation on tumor-relevant regions, reducing effective sequence length and improving both efficiency and performance. The soft masking weights patches by tumor probability, providing benefits of hard attention while maintaining differentiability.

\subsection{Multi-Task Learning in Medical Imaging}

Multi-task learning leverages shared representations across related tasks to improve generalization \cite{caruana1997multitask}. In medical imaging, several works explored joint segmentation and classification. H-DenseUNet \cite{li2018hdenseunet} used hybrid densely connected networks for liver lesion segmentation and classification. Y-Net \cite{mehta2018ynet} proposed dual-branch architectures with shared encoders.

However, most multi-task approaches use simple feature sharing through common encoders or independent task heads, without explicit mechanisms to guide one task using information from another. **AMT-UNet introduces ROI-guided multi-task learning** where segmentation predictions explicitly gate the classification input through multiplicative masking. This focuses grade prediction on tumor regions while avoiding background interference. Importantly, we employ gradient stopping ($\text{detach}$) to prevent classification gradients from affecting segmentation, eliminating task interference that often plagues multi-task learning. This one-way information flow (segmentation $\rightarrow$ classification) proves more effective than bidirectional coupling.

\subsection{Loss Functions for Medical Image Segmentation}

Cross-entropy loss, while standard in classification, struggles with severe class imbalance in segmentation. Dice loss \cite{milletari2016vnet} directly optimizes the Dice coefficient, providing better handling of imbalanced classes. Focal loss \cite{lin2017focal} down-weights easy examples through a modulating factor, forcing the model to focus on hard misclassified regions. Tversky loss \cite{salehi2017tversky} generalizes Dice by asymmetrically weighting false positives and false negatives. Boundary loss \cite{kervadec2019boundary} uses distance transforms to emphasize boundary precision.

Recent works combine multiple loss components. However, the challenge lies in selecting compatible losses and determining appropriate weights. **AMT-UNet employs a principled combination of four complementary losses**: (1) Dice loss provides stable region-based optimization as the foundation; (2) focal loss with strong focusing ($\gamma=3.0$, higher than standard $\gamma=2.0$) aggressively addresses class imbalance; (3) IoU loss weighted $2.5\times$ directly optimizes the evaluation metric, resolving the train-test objective mismatch; (4) boundary loss ensures precise edge delineation critical for clinical utility. Our weights are motivated by the specific failure modes of brain tumor segmentation: IoU is weighted heavily because it is stricter than Dice and the primary evaluation metric, while boundary loss receives moderate weight (0.6) to improve edge precision without overwhelming the region-based objectives.

\subsection{Deep Supervision and Multi-Scale Learning}

Deep supervision \cite{lee2015deeply} adds auxiliary loss functions at intermediate layers to improve gradient flow and accelerate convergence. U-Net++ \cite{zhou2018unet++} extensively uses deep supervision across nested decoder paths. For medical imaging, deep supervision has proven effective for training very deep networks and improving feature learning at multiple scales.

**AMT-UNet combines deep supervision with multi-scale feature fusion** in a unified framework. We attach auxiliary segmentation heads at three decoder levels ($64 \times 64$, $128 \times 128$, $256 \times 256$) with weight 0.3, ensuring consistent supervision across scales. Additionally, we fuse features from all four decoder levels before the final prediction, explicitly combining multi-resolution information. This dual approach—supervising at multiple scales and fusing multi-scale features—provides stronger gradient signals and richer representations than either technique alone. The fusion addresses the semantic gap between shallow (detail-rich) and deep (context-rich) features that pure skip connections fail to bridge.
